\begin{thebibliography}{99}

\bibitem{BB81} A. Beilinson and J. Bernstein, \emph{Localisation de $\mathfrak{g}$-modules}, C. R. Acad. Sci. Paris S\'er. I Math. \textbf{292} (1981), no. 1, 15--18.

\bibitem{BD04} A. Beilinson and V. Drinfeld, \emph{Chiral Algebras}, American Mathematical Society Colloquium Publications, vol. 51, American Mathematical Society, Providence, RI, 2004.

\bibitem{BGG76} I. N. Bernstein, I. M. Gel'fand, and S. I. Gel'fand, \emph{A certain category of $\mathfrak{g}$-modules}, Funktsional. Anal. i Prilozhen. \textbf{10} (1976), no. 2, 1--8.

\bibitem{BGS96} A. Beilinson, V. Ginzburg, and W. Soergel, \emph{Koszul duality patterns in representation theory}, J. Amer. Math. Soc. \textbf{9} (1996), no. 2, 473--527.

\bibitem{BK81} J.-L. Brylinski and M. Kashiwara, \emph{Kazhdan--Lusztig conjecture and holonomic systems}, Invent. Math. \textbf{64} (1981), no. 3, 387--410.

\bibitem{DGMS75} P. Deligne, P. Griffiths, J. Morgan, and D. Sullivan, \emph{Real homotopy theory of K\"ahler manifolds}, Invent. Math. \textbf{29} (1975), no. 3, 245--274.

\bibitem{BernhartRiordan} F. R. Bernhart, \emph{Catalan, Motzkin, and Riordan numbers}, Discrete Math. \textbf{204} (1999), no. 1--3, 73--112.

\bibitem{BW83} A. Björner and M. L. Wachs, On lexicographically shellable posets, \emph{Trans. Amer. Math. Soc.} \textbf{277} (1983), no. 1, 323--331.
 
\bibitem{FBZ04} E. Frenkel and D. Ben-Zvi, \emph{Vertex Algebras and Algebraic Curves}, Mathematical Surveys and Monographs, vol. 88, American Mathematical Society, Providence, RI, 2004.
 
\bibitem{FlajoletSedgewick} P. Flajolet and R. Sedgewick, \emph{Analytic Combinatorics}, Cambridge University Press, Cambridge, 2009.

\bibitem{FM94} W. Fulton and R. MacPherson, A compactification of configuration spaces, \emph{Ann. of Math.} (2) \textbf{139} (1994), no. 1, 183--225.
 

\bibitem{OS80} P. Orlik and L. Solomon, Combinatorics and topology of complements of hyperplanes, \emph{Invent. Math.} \textbf{56} (1980), no. 2, 167--189.
 
\bibitem{Sta97} R. P. Stanley, \emph{Enumerative Combinatorics}, vol. 1, Cambridge Studies in Advanced Mathematics, vol. 49, Cambridge University Press, Cambridge, 1997.

\bibitem{Sta63} J. D. Stasheff, \emph{Homotopy associativity of H-spaces. I, II}, Trans. Amer. Math. Soc. \textbf{108} (1963), 275--292, 293--312.

\bibitem{LV} J.-L. Loday and B. Vallette, \emph{Algebraic Operads}, Grundlehren der mathematischen Wissenschaften, vol. 346, Springer, 2012.

\bibitem{Loday98}
J.-L. Loday, \emph{Cyclic Homology}, second edition, Grundlehren der mathematischen Wissenschaften, vol. 301, Springer-Verlag, Berlin, 1998.

\bibitem{Connes85}
A. Connes, \emph{Noncommutative differential geometry}, Inst. Hautes \'Etudes Sci. Publ. Math. No. 62 (1985), 257--360.

\bibitem{Eyn04} B. Eynard, \emph{All genus correlation functions for the hermitian 1-matrix model}, J. High Energy Phys. (2004), no. 11, 031.

\bibitem{CEO06} L. Chekhov, B. Eynard, and N. Orantin, \emph{Free energy topological expansion for the 2-matrix model}, J. High Energy Phys. (2006), no. 12, 053.

\bibitem{EO07} B. Eynard and N. Orantin, \emph{Invariants of algebraic curves and topological expansion}, Commun. Number Theory Phys. \textbf{1} (2007), no. 2, 347--452.

\bibitem{GG11} M. R. Gaberdiel and R. Gopakumar, \emph{An $\mathrm{AdS}_3$ dual for minimal model CFTs}, Phys.\ Rev.\ D \textbf{83} (2011), 066007.

\bibitem{Get95} E. Getzler, \emph{Operads and moduli spaces of genus 0 Riemann surfaces}, in \emph{The Moduli Space of Curves} (Texel Island, 1994), Progr. Math. \textbf{129}, Birkhäuser, Boston, 1995, pp. 199--230.

\bibitem{GeK98} E. Getzler and M. M. Kapranov, \emph{Modular operads}, Compositio Math. \textbf{110} (1998), no. 1, 65--126.

\bibitem{GJ} E. Getzler and J.D.S. Jones, \emph{Operads, homotopy algebra and iterated integrals for double loop spaces}, arXiv:hep-th/9403055, 1994.

\bibitem{FG06} E. Frenkel and D. Gaitsgory, \emph{Local geometric Langlands correspondence and affine Kac--Moody algebras}, in Algebraic Geometry and Number Theory, Progr. Math. \textbf{253}, Birkh\"auser, Boston, 2006, 69--260.

\bibitem{FG12} J. Francis and D. Gaitsgory, \emph{Chiral Koszul duality}, Selecta Math. \textbf{18} (2012), no. 1, 27--87.

\bibitem{CG17} K. Costello and O. Gwilliam, \emph{Factorization algebras in quantum field theory}, Vols. 1-2, Cambridge University Press, 2017.

\bibitem{BK} A. A. Belavin and V. G. Knizhnik, \emph{Complex geometry and theory of quantum strings}, 
  Zh. Eksp. Teor. Fiz. 91 (1986), 364--390.
  
\bibitem{Verlinde} E. Verlinde, \emph{Fusion rules and modular transformations in 2D conformal field theory}, 
  Nuclear Phys. B 300 (1988), 360--376.

\bibitem{Zhu} Y. Zhu, \emph{Modular invariance of characters of vertex operator algebras},
  J. Amer. Math. Soc. 9 (1996), 237--302.

\bibitem{FZ92} I. Frenkel and Y. Zhu, \emph{Vertex operator algebras associated to representations of affine and Virasoro algebras}, Duke Math. J. \textbf{66} (1992), no. 1, 123--168.

\bibitem{GK94} V. Ginzburg and M. Kapranov, \emph{Koszul duality for operads}, Duke Math. J. \textbf{76} (1994), no. 1, 203--272.

\bibitem{HTT08} R. Hotta, K. Takeuchi, and T. Tanisaki, \emph{$D$-Modules, Perverse Sheaves, and Representation Theory}, Progress in Mathematics, vol. 236, Birkhäuser, Boston, MA, 2008.

\bibitem{Kon94} M. Kontsevich, \emph{Feynman diagrams and low-dimensional topology}, First European Congress of Mathematics, Vol. II (Paris, 1992), 97--121, Progr. Math., 120, Birkhäuser, Basel, 1994.

\bibitem{Kon99} M. Kontsevich, \emph{Operads and motives in deformation quantization}, Lett. Math. Phys. \textbf{48} (1999), no. 1, 35--72.

\bibitem{Kon03} M. Kontsevich, \emph{Deformation quantization of Poisson manifolds}, Lett. Math. Phys. \textbf{66} (2003), no. 3, 157--216.

\bibitem{Ara07} T. Arakawa, \emph{Representation theory of W-algebras}, Invent. Math. \textbf{169} (2007), no. 2, 219--320.

\bibitem{Ara15} T. Arakawa, \emph{Rationality of W-algebras: principal nilpotent cases}, Ann. Math. \textbf{182} (2015), no. 2, 565--694.

\bibitem{May72} J. P. May, \emph{The geometry of iterated loop spaces}, Lectures Notes in Mathematics, Vol. 271. Springer-Verlag, Berlin-New York, 1972.

\bibitem{BV73} J. M. Boardman and R. M. Vogt, \emph{Homotopy invariant algebraic structures on topological spaces}, Lecture Notes in Mathematics, Vol. 347. Springer-Verlag, Berlin-New York, 1973.

\bibitem{Coh76} F. R. Cohen, \emph{The homology of $C_{n+1}$-spaces, $n \geq 0$}, Lecture Notes in Math., Vol. 533, Springer-Verlag, Berlin-Heidelberg-New York, 1976, pp. 207--351.

\bibitem{Klingen67} H. Klingen, \emph{Zum Darstellungssatz f\"ur Siegelsche Modulformen}, Math.\ Z.\ \textbf{102} (1967), 30--43.

\bibitem{HA} J. Lurie, \emph{Higher Algebra}, available at http://www.math.harvard.edu/~lurie/.

\bibitem{GR17}
Dennis Gaitsgory and Nick Rozenblyum.
\textit{A Study in Derived Algebraic Geometry}, Mathematical Surveys and Monographs, vols. 221.1--221.2, American Mathematical Society, 2017.

\bibitem{AF15} D. Ayala and J. Francis, \emph{Factorization homology of topological manifolds}, J. Topology \textbf{8} (2015), no. 4, 1045--1084.

\bibitem{BPZ84} A. A. Belavin, A. M. Polyakov and A. B. Zamolodchikov, \emph{Infinite conformal symmetry in two-dimensional quantum field theory}, Nuclear Phys. B \textbf{241} (1984), no. 2, 333--380.

\bibitem{FP00} C. Faber and R. Pandharipande, \emph{Hodge integrals and Gromov--Witten theory}, Invent. Math. \textbf{139} (2000), no. 1, 173--199.

\bibitem{FF} B. Feigin and E. Frenkel, \emph{Quantization of the Drinfeld--Sokolov reduction}, Phys. Lett. B \textbf{246} (1990), 75--81.

\bibitem{Drinfeld85} V. Drinfeld, \emph{Hopf algebras and the quantum Yang--Baxter equation}, Soviet Math. Dokl. \textbf{32} (1985), 254--258.

\bibitem{DS} V. Drinfeld and V. Sokolov, \emph{Lie algebras and equations of Korteweg-de Vries type}, Soviet Math. Dokl. \textbf{23} (1981), 457--462.

\bibitem{Koszul50} J.-L. Koszul, \emph{Homologie et cohomologie des alg\`ebres de Lie}, Bull. Soc. Math. France \textbf{78} (1950), 65--127.

\bibitem{Priddy70} S. Priddy, \emph{Koszul resolutions}, Trans. Amer. Math. Soc. \textbf{152} (1970), 39--60.

\bibitem{KRW} V. Kac, S.-S. Roan, and M. Wakimoto, \emph{Quantum reduction for affine superalgebras}, Comm. Math. Phys. \textbf{241} (2003), 307--342.

\bibitem{Positselski} L. Positselski, \emph{Two kinds of derived categories, Koszul duality, and comodule-contramodule correspondence}, Mem. Amer. Math. Soc. \textbf{212} (2011), no. 996.

% Feynman Diagrams and Worldline Formalism

\bibitem{worldline-formalism}
C. Schubert, \emph{Perturbative quantum field theory in the string-inspired formalism}, 
Phys. Rept. \textbf{355} (2001) 73-234, arXiv:hep-th/0101036.

% BV-BRST Formalism

\bibitem{batalin-vilkovisky}
I. A. Batalin and G. A. Vilkovisky, \emph{Gauge algebra and quantization}, 
Phys. Lett. B \textbf{102} (1981) 27-31.

\bibitem{gwilliam-thesis}
O. Gwilliam, \emph{Factorization algebras and free field theories}, 
PhD thesis, Northwestern University (2012).

\bibitem{mnev-bv}
P. Mnev, \emph{Lectures on Batalin--Vilkovisky formalism and its applications 
in topological quantum field theory}, arXiv:1707.08096.

% Holomorphic-Topological Field Theory

\bibitem{costello-gaiotto}
K. Costello and D. Gaiotto, \emph{Twisted holography}, 
arXiv:1812.09257.

\bibitem{gaiotto-rapchak}
D. Gaiotto and J. Rapchak, \emph{Vertex Algebras at the Corner}, 
JHEP \textbf{01} (2019) 160, arXiv:1703.00982.

\bibitem{costello-yamazaki}
K. Costello and M. Yamazaki, \emph{Gauge Theory And Integrability, III}, 
arXiv:1908.02289.

\bibitem{oh-yagi}
J. Oh and Y. Yagi, \emph{Chiral algebra, localization, modularity, surface defects, 
and all that}, arXiv:1910.11261.

% AGT and W-Algebras

\bibitem{arakawa-lectures}
T. Arakawa, \emph{Introduction to W-algebras and their representation theory}, 
in Perspectives in Lie Theory, Springer INdAM Series \textbf{19} (2017) 179-250, 
arXiv:1605.00138.

\bibitem{arakawa-higgs}
T. Arakawa and A. Molev, \emph{Explicit generators in rectangular affine 
$\mathcal{W}$-algebras of type A}, 
Lett. Math. Phys. \textbf{107} (2017) 47-59, arXiv:1403.1017.

\bibitem{beem-et-al}
C. Beem, M. Lemos, P. Liendo, W. Peelaers, L. Rastelli, and B. C. van Rees, 
\emph{Infinite chiral symmetry in four dimensions}, 
Comm. Math. Phys. \textbf{336} (2015) 1359-1433, arXiv:1312.5344.

% Yangians and Quantum Groups

\bibitem{molev-yangians}
A. Molev, \emph{Yangians and Classical Lie Algebras}, 
Mathematical Surveys and Monographs \textbf{143}, AMS (2007).

\bibitem{nakajima-quiver}
H. Nakajima, \emph{Quiver varieties and finite dimensional representations of 
quantum affine algebras}, 
J. Amer. Math. Soc. \textbf{14} (2001) 145-238, arXiv:math/9912158.

% Additional Key References

\bibitem{EGNO}
P. Etingof, S. Gelaki, D. Nikshych, and V. Ostrik,
\emph{Tensor Categories}, Mathematical Surveys and Monographs,
vol.~205, American Mathematical Society, 2015.

\bibitem{etingof-kazhdan}
P. Etingof and D. Kazhdan, \emph{Quantization of Lie bialgebras, V},
Selecta Math. (N.S.) \textbf{6} (2000) 105-130, arXiv:math/9808121.

% Phase 3 Additional References

\bibitem{Kac}
V. G. Kac, \emph{Infinite Dimensional Lie Algebras},
3rd edition, Cambridge University Press, 1990.

\bibitem{Kal08}
D. Kaledin, \emph{Non-commutative Hodge-to-de Rham degeneration via the method of Deligne--Illusie},
Pure Appl. Math. Q. \textbf{4} (2008), no. 3, 785--875.

\bibitem{Zamolodchikov}
A. B. Zamolodchikov, \emph{Infinite additional symmetries in two-dimensional conformal quantum field theory},
Theor. Math. Phys. \textbf{65} (1985) 1205-1213.

\bibitem{Bouwknegt-Schoutens}
P. Bouwknegt and K. Schoutens, \emph{W-symmetry in conformal field theory},
Phys. Rep. \textbf{223} (1993) 183-276, arXiv:hep-th/9210010.

\bibitem{Feigin-Frenkel}
B. Feigin and E. Frenkel, \emph{Affine Kac--Moody algebras at the critical level and Gelfand--Dikii algebras},
Int. J. Mod. Phys. A \textbf{7} (1992) Suppl. 1A, 197-215.

\bibitem{LodayVallette} J.-L. Loday and B. Vallette,
\emph{Algebraic Operads},
Grundlehren der mathematischen Wissenschaften 346, Springer, 2012.

\bibitem{Hormander} L. Hörmander,
\emph{The Analysis of Linear Partial Differential Operators I: Distribution Theory and Fourier Analysis},
Classics in Mathematics, Springer, 2003.
See especially Chapter 8 on multiplication of distributions.

\bibitem{Mel93}
Richard B. Melrose.
\textit{The Atiyah--Patodi--Singer index theorem}.
A K Peters, 1993.

\bibitem{Sch66}
Laurent Schwartz.
\textit{Théorie des distributions}.
Hermann, 1966.

\bibitem{KS94}
Masaki Kashiwara and Pierre Schapira.
\textit{Sheaves on Manifolds}.
Grundlehren der Mathematischen Wissenschaften 292.
Springer-Verlag, 1994.

\bibitem{FZ87}
V. A. Fateev and A. B. Zamolodchikov.
\textit{Conformal quantum field theory models in two dimensions having $Z_3$ symmetry}.
Nuclear Physics B, 280:644-660, 1987.

\bibitem{Zhu96}
Yongchang Zhu.
\textit{Modular invariance of characters of vertex operator algebras}.
J. Amer. Math. Soc., 9(1):237-302, 1996.

\bibitem{KR87}
Victor G. Kac and Ashok K. Raina.
\textit{Bombay Lectures on Highest Weight Representations of Infinite Dimensional Lie Algebras}.
Advanced Series in Mathematical Physics 2.
World Scientific, 1987.

% ================================================================
% PATCH 013-015: NEW BIBLIOGRAPHY ENTRIES
% ================================================================

\bibitem{GLZ22}
Z. Gui, S. Li, and K. Zeng, \emph{Quadratic duality for chiral algebras},
Advances in Mathematics \textbf{451} (2024), 109779, arXiv:2212.11252.

\bibitem{LurieHTT}
Jacob Lurie.
\textit{Higher Topos Theory}.
Annals of Mathematics Studies 170, Princeton University Press, 2009.
Note: $(\infty,1)$-categories and homotopy coherent structures.

\bibitem{LurieHA}
Jacob Lurie.
\textit{Higher Algebra}.
Available at \url{https://www.math.ias.edu/~lurie/papers/HA.pdf}, 2017.
Note: $E_n$-algebras, Koszul duality, and stable $\infty$-categories.

% Positselski-curved: removed (duplicate of Positselski11)

\bibitem{Fresse-operads}
Benoit Fresse.
\textit{Modules over Operads and Functors}.
Springer, 2009.
Note: Comprehensive treatment of operadic structures.

\bibitem{CG-vol2}
Kevin Costello and Owen Gwilliam.
\textit{Factorization Algebras in Quantum Field Theory, Volume 2}.
Cambridge University Press, 2021.
Note: BV formalism and quantum corrections.

% ================================================================
% PATCH 017: ARNOLD RELATIONS HISTORICAL SOURCES
% ================================================================

\bibitem{Arnold69}
V. I. Arnold.
\textit{The cohomology ring of the colored braid group}.
Mat. Zametki, 5:227-231, 1969.
Note: English translation: Math. Notes 5 (1969), 138-140. 
Original discovery of the three-term relations for configuration spaces.

\bibitem{Brieskorn73}
Egbert Brieskorn.
\textit{Sur les groupes de tresses [d'après V. I. Arnold]}.
Séminaire Bourbaki, Lecture Notes in Mathematics 317:21-44, Springer, 1973.
Note: Generalization to hyperplane arrangements, nine-term exact sequence, 
singularity theory connections. Seminal work connecting Arnold's relations 
to broader arrangement theory.

\bibitem{OrlikTerao92}
Peter Orlik and Hiroaki Terao.
\newblock \textit{Arrangements of Hyperplanes}.
\newblock Grundlehren der mathematischen Wissenschaften, Vol.~300.
\newblock Springer-Verlag, 1992.
\newblock ISBN 978-3-540-55259-9.

\bibitem{GoMa88}
Mark Goresky and Robert MacPherson.
\textit{Stratified Morse theory}.
Ergebnisse der Mathematik und ihrer Grenzgebiete 14, Springer-Verlag, 1988.
Note: Stratified spaces perspective on arrangement complements, 
intersection cohomology approach to Arnold relations.

\bibitem{Kontsevich97}
Maxim Kontsevich.
\textit{Formality conjecture}.
In: Deformation theory and symplectic geometry, Math. Phys. Stud. 20:139-156, Kluwer, 1997.
Note: Configuration space integrals, formality theorem. 
Arnold relations ensure integrals satisfy $d^2 = 0$.

\bibitem{BaKi01}
B. Bakalov and A. Kirillov, Jr.
\textit{Lectures on Tensor Categories and Modular Functors}.
University Lecture Series \textbf{21}, American Mathematical Society,
Providence, RI, 2001.

\bibitem{Cohen76}
Frederick R. Cohen.
\textit{The homology of $C_{n+1}$-spaces, $n \geq 0$}.
In: The homology of iterated loop spaces, Lecture Notes in Mathematics 533:207-351, Springer, 1976.
Note: Homology of configuration spaces, connection to iterated loop spaces. 
Foundational for understanding topology of $\text{Conf}_n$.

\bibitem{Totaro96}
Burt Totaro.
\textit{Configuration spaces of algebraic varieties}.
Topology, 35(4):1057-1067, 1996.
Note: Configuration spaces over arbitrary varieties, 
extension beyond $\mathbb{C}$. Relevant for our higher genus work.

% FultonMacPherson94 consolidated into FM94 above

% ================================================================
% PATCH 017: HOLOMORPHIC-TOPOLOGICAL TWIST REFERENCES
% ================================================================

\bibitem{CL16}
Kevin Costello and Si Li.
\textit{Twisted supergravity and its quantization}.
arXiv:1606.00365, 2016.

\bibitem{Gai19}
Davide Gaiotto.
\textit{Twisted holography and vertex operator algebras at corners}.
arXiv:1903.00382, 2019.

\bibitem{PW22}
Natalie M. Paquette and Brian R. Williams.
\textit{On the definition of vertex algebras in holomorphic-topological twist}.
Communications in Mathematical Physics, 391:1185-1235, 2022.

\bibitem{ACL19}
Tomoyuki Arakawa, Thomas Creutzig, and Andrew R. Linshaw.
\textit{W-algebras as coset vertex algebras}.
Inventiones mathematicae, 218(1):145-195, 2019.

\bibitem{AGT09}
Luis F. Alday, Davide Gaiotto, and Yuji Tachikawa.
\textit{Liouville correlation functions from four-dimensional gauge theories}.
Letters in Mathematical Physics, 91(2):167-197, 2010.

\bibitem{FL88}
V. A. Fateev and S. L. Lukyanov.
\textit{The models of two-dimensional conformal quantum field theory with $Z_n$ symmetry}.
Int. J. Mod. Phys. A, 3:507-520, 1988.

\bibitem{Mizoguchi89}
S. Mizoguchi.
\textit{Determinant formula and unitarity for the $W_3$ algebra}.
Phys. Lett. B, 222(2):226--230, 1989.

\bibitem{Weibel94}
Charles A. Weibel.
\textit{An Introduction to Homological Algebra}.
Cambridge Studies in Advanced Mathematics 38, Cambridge University Press, 1994.

% ==============================================================================
% ARAKAWA'S WORK: COMPLETE BIBLIOGRAPHY (PATCH 051)
% ==============================================================================

\bibitem{Arakawa2005Duke}
Tomoyuki Arakawa.
\textit{Representation theory of superconformal algebras and the Kac--Roan--Wakimoto conjecture}.
Duke Mathematical Journal, 130(3):435-478, 2005.

\bibitem{Arakawa2007Invent}
Tomoyuki Arakawa.
\textit{Representation theory of $W$-algebras}.
Inventiones Mathematicae, 169(2):219-320, 2007.

\bibitem{Arakawa2015ICM}
Tomoyuki Arakawa.
\textit{Associated varieties of modules over Kac--Moody algebras and $C_2$-cofiniteness of $W$-algebras}.
In: Proceedings of the International Congress of Mathematicians---Seoul 2014, Vol. II, pages 1109-1125, 2015.

\bibitem{Arakawa2016RationalAdmissible}
Tomoyuki Arakawa.
\textit{Rationality of admissible affine vertex algebras in the category $\mathcal{O}$}.
Duke Mathematical Journal, 165(1):67-93, 2016.

\bibitem{ArakawaCreutzigLinshaw2019}
Tomoyuki Arakawa, Thomas Creutzig, and Andrew R. Linshaw.
\textit{$W$-algebras as coset vertex algebras}.
Inventiones Mathematicae, 218:145-195, 2019.

\bibitem{AdamovicMilas2008}
Dra\v{z}en Adamovi\'c and Antun Milas.
\textit{On the triplet vertex algebra $\mathcal{W}(p)$}.
Advances in Mathematics, 217(6):2664--2699, 2008.

\bibitem{CreutzigRidout2013}
Thomas Creutzig and David Ridout.
\textit{Logarithmic conformal field theory: beyond an introduction}.
Journal of Physics A: Mathematical and Theoretical,
46(49):494006, 2013.

\bibitem{BCOV93}
M. Bershadsky, S. Cecotti, H. Ooguri, and C. Vafa.
\textit{Kodaira--Spencer theory of gravity and exact results for quantum string amplitudes}.
Communications in Mathematical Physics, 165(2):311-427, 1994.

\bibitem{Frenkel-Kac-Wakimoto92}
I. B. Frenkel, V. G. Kac, and M. Wakimoto.
\textit{Characters and fusion rules for $W$-algebras via quantized Drinfeld--Sokolov reduction}.
Communications in Mathematical Physics, 147(2):295-328, 1992.

\bibitem{Wakimoto86}
M. Wakimoto.
\textit{Fock representations of the affine Lie algebra $A_1^{(1)}$}.
Communications in Mathematical Physics, 104(4):605-609, 1986.

\bibitem{Kapustin-Witten06}
A. Kapustin and E. Witten.
\textit{Electric-magnetic duality and the geometric Langlands program}.
Communications in Number Theory and Physics, 1(1):1-236, 2007.

\bibitem{Verlinde88}
Erik Verlinde.
\textit{Fusion rules and modular transformations in 2D conformal field theory}.
Nuclear Physics B, 300:360-376, 1988.

\bibitem{MooreSeiberg89}
Gregory Moore and Nathan Seiberg.
\textit{Classical and quantum conformal field theory}.
Communications in Mathematical Physics, 123(2):177-254, 1989.

\bibitem{Costello-1705.02500v1}
Kevin Costello.
\textit{Holography and Koszul duality: the example of the M2 brane}.
arXiv preprint arXiv:1705.02500, 2017.
Note: Page 7: ``The Koszul dual of $C^*(g)$ is the universal enveloping algebra $U(g)$''.

\bibitem{Voronov99}
Alexander Voronov.
\textit{The Swiss-cheese operad}.
In: Homotopy invariant algebraic structures, Contemp.\ Math.\ 239, AMS, 1999, pp.~365--373.
Note: Introduces the Swiss-cheese operad SC$_n$ governing pairs (E$_n$-algebra, module).

\bibitem{Wei99}
Michael Weiss.
\textit{Embeddings from the point of view of immersion theory, Part I}.
Geometry \& Topology, 3:67--101, 1999.
Note: Introduces the Weiss topology (covering families where every finite subset is contained in a single open) for manifold calculus of functors.

% ------------------------------------------------------------------
% Legacy citation-key aliases (Phase 0 hygiene)
% These preserve old in-text keys while keeping a single source entry.
% ------------------------------------------------------------------
\bibitem{AF} D. Ayala and J. Francis, \emph{Factorization homology of topological manifolds}, J. Topology \textbf{8} (2015), no. 4, 1045--1084.

\bibitem{Ara12} T. Arakawa, \emph{Introduction to $W$-algebras and their representation theory}, in Perspectives in Lie Theory, Springer INdAM Series \textbf{19} (2017), 179--250, arXiv:1605.00138.

\bibitem{Arakawa} T. Arakawa, \emph{Chiral algebras of class $\mathcal{S}$ and Moore--Tachikawa varieties}, arXiv:1811.01577, 2018.

\bibitem{Arakawa-W-algebras} T. Arakawa, \emph{Introduction to W-algebras and their representation theory}, in Perspectives in Lie Theory, Springer INdAM Series \textbf{19} (2017), 179--250, arXiv:1605.00138.

\bibitem{Arakawa-lectures} T. Arakawa, \emph{Introduction to W-algebras and their representation theory}, in Perspectives in Lie Theory, Springer INdAM Series \textbf{19} (2017), 179--250, arXiv:1605.00138.

\bibitem{Arakawa17} T. Arakawa, \emph{Chiral algebras of class $\mathcal{S}$ and Moore--Tachikawa varieties}, arXiv:1811.01577, 2018.

\bibitem{Arn69} V. I. Arnold, \emph{The cohomology ring of the colored braid group}, Mat. Zametki \textbf{5} (1969), 227--231.

\bibitem{BD} A. Beilinson and V. Drinfeld, \emph{Chiral Algebras}, American Mathematical Society Colloquium Publications, vol. 51, American Mathematical Society, Providence, RI, 2004.

\bibitem{FF-wakimoto} I. B. Frenkel, V. G. Kac, and M. Wakimoto, \emph{Characters and fusion rules for $W$-algebras via quantized Drinfeld--Sokolov reduction}, Comm. Math. Phys. \textbf{147} (1992), no. 2, 295--328.

\bibitem{FG} J. Francis and D. Gaitsgory, \emph{Chiral Koszul duality}, Selecta Math. \textbf{18} (2012), no. 1, 27--87.

\bibitem{HTT} J. Lurie, \emph{Higher Topos Theory}, Annals of Mathematics Studies, vol. 170, Princeton University Press, 2009.

\bibitem{KS90} M. Kashiwara and P. Schapira, \emph{Sheaves on Manifolds}, Grundlehren der Mathematischen Wissenschaften, vol. 292, Springer-Verlag, 1990.

\bibitem{Kontsevich03} M. Kontsevich, \emph{Deformation quantization of Poisson manifolds}, Lett. Math. Phys. \textbf{66} (2003), no. 3, 157--216.

\bibitem{Li96}
H. Li, \emph{Local systems of vertex operators, vertex superalgebras and modules}, J. Pure Appl. Algebra \textbf{109} (1996), no. 2, 143--195.

\bibitem{LV12} J.-L. Loday and B. Vallette, \emph{Algebraic Operads}, Grundlehren der mathematischen Wissenschaften, vol. 346, Springer, 2012.

\bibitem{fay1973theta}
John D. Fay.
\textit{Theta Functions on Riemann Surfaces}.
Lecture Notes in Mathematics, vol. 352, Springer-Verlag, Berlin-New York, 1973.

\bibitem{Fay73}
John D. Fay.
\textit{Theta Functions on Riemann Surfaces}.
Lecture Notes in Mathematics, vol. 352, Springer-Verlag, Berlin-New York, 1973.

\bibitem{Igusa62}
Jun-ichi Igusa.
\textit{On Siegel modular forms of genus two}.
American Journal of Mathematics \textbf{84} (1962), 175--200.

\bibitem{Igusa72}
Jun-ichi Igusa.
\textit{Theta Functions}.
Grundlehren der mathematischen Wissenschaften, vol. 194, Springer-Verlag, Berlin-New York, 1972.

\bibitem{Freitag83}
Eberhard Freitag.
\textit{Siegel's Modular Forms}.
Grundlehren der mathematischen Wissenschaften, vol. 254, Springer-Verlag, Berlin, 1983.

% ------------------------------------------------------------------
% Legacy keys with verified metadata
% ------------------------------------------------------------------

\bibitem{CP2020}
Kevin Costello and Natalie M. Paquette.
\textit{Twisted Supergravity and Koszul Duality: A Case Study in AdS$_3$}.
Communications in Mathematical Physics \textbf{384} (2021), 279--339.
arXiv:2001.02177. DOI: 10.1007/s00220-021-04065-3.

\bibitem{DM69}
Pierre Deligne and David Mumford.
\textit{The irreducibility of the space of curves of given genus}.
Publications Mathématiques de l'IHÉS \textbf{36} (1969), 75--109.

\bibitem{DeligneM69} P. Deligne and D. Mumford, \emph{The irreducibility of the space of curves of given genus}, Publ. Math. IH\'ES \textbf{36} (1969), 75--109.

\bibitem{Knudsen83} F. F. Knudsen, \emph{The projectivity of the moduli space of stable curves, II: The stacks $M_{g,n}$}, Math. Scand. \textbf{52} (1983), 161--199.

\bibitem{DMVV}
Robbert Dijkgraaf, Gregory Moore, Erik Verlinde, and Herman Verlinde.
\textit{Elliptic genera of symmetric products and second quantized strings}.
Communications in Mathematical Physics \textbf{185} (1997), 197--209.
arXiv:hep-th/9608096. DOI: 10.1007/s002200050087.

\bibitem{Deligne-Illusie}
Pierre Deligne and Luc Illusie.
\textit{Relèvements modulo $p^2$ et décomposition du complexe de de Rham}.
Inventiones Mathematicae \textbf{89} (1987), 247--270.

\bibitem{FF92} B. Feigin and E. Frenkel, \emph{Affine Kac--Moody algebras at the critical level and Gelfand--Dikii algebras}, Int. J. Mod. Phys. A \textbf{7} (1992) Suppl. 1A, 197--215.

\bibitem{Ger63}
Murray Gerstenhaber.
\textit{The cohomology structure of an associative ring}.
Annals of Mathematics (2) \textbf{78} (1963), 267--288.
DOI: 10.2307/1970343.

\bibitem{Keller01}
Bernhard Keller.
\textit{Introduction to $A$-infinity algebras and modules}.
Homology, Homotopy and Applications \textbf{3} (2001), no.~1, 1--35.
arXiv:math/9910179.

\bibitem{Merkulov99}
Sergei Merkulov.
\textit{Strong homotopy algebras of a K\"ahler manifold}.
International Mathematics Research Notices \textbf{1999}, no.~3, 153--164.
arXiv:math/9809172.

\bibitem{Mumford83}
David Mumford.
\textit{Towards an enumerative geometry of the moduli space of curves}.
In \textit{Arithmetic and Geometry, Vol. II}, Progress in Mathematics \textbf{36},
Birkhäuser, Boston, MA, 1983, 271--328.

\bibitem{Mumford84}
David Mumford.
\textit{Tata Lectures on Theta II: Jacobian Theta Functions and Differential Equations}.
Progress in Mathematics \textbf{43}, Birkhäuser Boston, 1984.

\bibitem{QuadDual} Z. Gui, S. Li, and K. Zeng, \emph{Quadratic duality for chiral algebras}, Advances in Mathematics \textbf{451} (2024), 109779, arXiv:2212.11252.

\bibitem{Sei88}
Nathan Seiberg.
\textit{Observations on the moduli space of superconformal field theories}.
Nuclear Physics B \textbf{303} (1988), 286--304.
DOI: 10.1016/0550-3213(88)90183-6.

\bibitem{SGA4}
M. Artin, A. Grothendieck, and J.-L. Verdier.
\textit{Théorie des topos et cohomologie étale des schémas (SGA 4), Tome 3:
Exposés IX à XIX}.
Lecture Notes in Mathematics, vol. 305, Springer-Verlag, Berlin-New York, 1973.
 
\bibitem{Watts-W4}
Peter Bowcock and Gerard M. T. Watts.
\textit{On the classification of quantum $W$-algebras}.
Nuclear Physics B \textbf{379} (1992), 63--95.
DOI: 10.1016/0550-3213(92)90590-8.

\bibitem{BFN18} A. Braverman, M. Finkelberg, and H. Nakajima, \emph{Towards a mathematical definition of Coulomb branches of $3$-dimensional $\mathcal{N}=4$ gauge theories, II}, Adv. Theor. Math. Phys. \textbf{22} (2018), no. 5, 1071--1147, arXiv:1601.03586.

\bibitem{Finkelberg96} M. Finkelberg, \emph{An equivalence of fusion categories}, Geom. Funct. Anal. \textbf{6} (1996), no. 2, 249--267.

\bibitem{AF20} D. Ayala and J. Francis, \emph{A factorization homology primer}, Handbook of Homotopy Theory, CRC Press, 2020, 39--101, arXiv:1903.10961.

\bibitem{FF84} B. L. Feigin and D. B. Fuchs, \emph{Verma modules over the Virasoro algebra}, in Topology (Leningrad, 1982), Lecture Notes in Math. \textbf{1060}, Springer, Berlin, 1984, 230--245.

\bibitem{KW88} V. G. Kac and M. Wakimoto, \emph{Modular invariant representations of infinite-dimensional Lie algebras and superalgebras}, Proc. Nat. Acad. Sci. USA \textbf{85} (1988), no. 14, 4956--4960.

\bibitem{Kee92} S. Keel, \emph{Intersection theory of moduli space of stable $n$-pointed curves of genus zero}, Trans. Amer. Math. Soc. \textbf{330} (1992), no. 2, 545--574.

\bibitem{GT72} P. Goddard and C. B. Thorn, \emph{Compatibility of the dual pomeron with unitarity and the absence of ghosts in the dual resonance model}, Phys. Lett. B \textbf{40} (1972), 235--238.

\bibitem{Pol98} J. Polchinski, \emph{String Theory}, Vols. 1--2, Cambridge University Press, Cambridge, 1998.

% Pos11: removed (duplicate of Positselski11)

\bibitem{Francis2013} J. Francis, \emph{The tangent complex and Hochschild cohomology of $E_n$-rings}, Compos. Math. \textbf{149} (2013), no. 3, 430--480.

\bibitem{FK80} I. B. Frenkel and V. G. Kac, \emph{Basic representations of affine Lie algebras and dual resonance models}, Invent. Math. \textbf{62} (1980), no. 1, 23--66.

\bibitem{Se81} G. Segal, \emph{Unitary representations of some infinite-dimensional groups}, Comm. Math. Phys. \textbf{80} (1981), no. 3, 301--342.

\bibitem{HLZ} Y.-Z. Huang, J. Lepowsky, and L. Zhang, \emph{Logarithmic tensor category theory for generalized modules for a conformal vertex algebra}, Parts I--VIII, preprint series, arXiv:1012.4193 et seq.

\bibitem{costello-renormalization} K. Costello, \emph{Renormalization and Effective Field Theory}, Mathematical Surveys and Monographs \textbf{170}, American Mathematical Society, 2011.

\bibitem{Positselski11} L. Positselski, \emph{Two kinds of derived categories, Koszul duality, and comodule-contramodule correspondence}, Mem. Amer. Math. Soc. \textbf{212} (2011), no. 996.

\bibitem{Boardman-conditional} J. M. Boardman, \emph{Conditionally convergent spectral sequences}, in \emph{Homotopy Invariant Algebraic Structures} (Baltimore, MD, 1998), Contemp. Math. \textbf{239}, Amer. Math. Soc., Providence, RI, 1999, pp. 49--84.

% Duplicates of Costello2017, SGA4, LV12 removed (already defined above as aliases)

% ================================================================
% SESSION 27: RESOLVE ALL UNDEFINED CITATIONS
% ================================================================

\bibitem{AY10}
H. Awata and Y. Yamada.
\textit{Five-dimensional AGT conjecture and the deformed Virasoro algebra}.
Journal of High Energy Physics \textbf{2010}, 125 (2010).
arXiv:0910.4431.

\bibitem{BZFN10}
D. Ben-Zvi, J. Francis, and D. Nadler.
\textit{Integral transforms and Drinfeld centers in derived algebraic geometry}.
Journal of the American Mathematical Society \textbf{23} (2010), no.~4, 909--966.
arXiv:0805.0157.

\bibitem{CWY18}
K. Costello, E. Witten, and M. Yamazaki.
\textit{Gauge theory and integrability, I}.
In \textit{Integrability, Quantization, and Geometry: I. Integrable Systems},
Proc. Sympos. Pure Math. \textbf{103.1}, AMS, 2021, 1--106.
arXiv:1709.09993.

\bibitem{FLM88}
I. Frenkel, J. Lepowsky, and A. Meurman.
\textit{Vertex Operator Algebras and the Monster}.
Pure and Applied Mathematics \textbf{134},
Academic Press, 1988.

\bibitem{Fre07}
E. Frenkel.
\textit{Langlands Correspondence for Loop Groups}.
Cambridge Studies in Advanced Mathematics \textbf{103},
Cambridge University Press, 2007.

\bibitem{FT87}
B. L. Feigin and B. L. Tsygan.
\textit{Additive K-theory}.
In \textit{K-theory, Arithmetic and Geometry} (Moscow, 1984--1986),
Lecture Notes in Mathematics \textbf{1289}, Springer, 1987, 67--209.

\bibitem{GW09}
D. Gaiotto and E. Witten.
\textit{S-duality of boundary conditions in $\mathcal{N}=4$ super Yang--Mills theory}.
Advances in Theoretical and Mathematical Physics \textbf{13} (2009), no.~3, 721--896.
arXiv:0807.3720.

\bibitem{Kadeishvili80} T. Kadeishvili, \emph{On the theory of homology of fiber spaces}, Uspekhi Mat. Nauk \textbf{35} (1980), no. 3, 183--188. English translation in Russian Math. Surveys \textbf{35} (1980), no. 3, 231--238.

\bibitem{KL79} D. Kazhdan and G. Lusztig, \emph{Representations of Coxeter groups and Hecke algebras}, Invent. Math. \textbf{53} (1979), no. 2, 165--184.

\bibitem{Drinfeld90}
V. G. Drinfeld, \emph{Quasi-Hopf algebras},
Leningrad Math. J. \textbf{1} (1990), no. 6, 1419--1457.

\bibitem{Kohno87}
T. Kohno, \emph{Monodromy representations of braid groups and Yang--Baxter equations},
Ann. Inst. Fourier (Grenoble) \textbf{37} (1987), no. 4, 139--160.

\bibitem{KL93}
D. Kazhdan and G. Lusztig.
\textit{Tensor structures arising from affine Lie algebras, I--IV}.
Journal of the American Mathematical Society \textbf{6} (1993), 905--947, 949--1011;
\textbf{7} (1994), 335--381, 383--453.

\bibitem{KT95} M. Kashiwara and T. Tanisaki, \emph{Kazhdan--Lusztig conjecture for affine Lie algebras with negative level}, Duke Math. J. \textbf{77} (1995), no. 1, 21--62.

\bibitem{Miyamoto04}
M. Miyamoto, \emph{Modular invariance of vertex operator algebras satisfying $C_2$-cofiniteness}, Duke Math. J. \textbf{122} (2004), no. 1, 51--91.

\bibitem{MO19}
D. Maulik and A. Okounkov.
\textit{Quantum Groups and Quantum Cohomology}.
Ast\'erisque \textbf{408} (2019).
arXiv:1211.1287.

\bibitem{RT91}
N. Yu. Reshetikhin and V. G. Turaev.
\textit{Invariants of $3$-manifolds via link polynomials and quantum groups}.
Inventiones Mathematicae \textbf{103} (1991), no.~3, 547--597.

\bibitem{SV13}
O. Schiffmann and E. Vasserot.
\textit{Cherednik algebras, W-algebras and the equivariant cohomology of the moduli space of instantons on $\mathbb{A}^2$}.
Publications Math\'ematiques de l'IH\'ES \textbf{118} (2013), 213--342.
arXiv:1202.2756.

\bibitem{SW99}
N. Seiberg and E. Witten.
\textit{String theory and noncommutative geometry}.
Journal of High Energy Physics \textbf{1999}, no.~9, 032.
arXiv:hep-th/9908142.

\bibitem{Wit89}
E. Witten.
\textit{Quantum field theory and the Jones polynomial}.
Communications in Mathematical Physics \textbf{121} (1989), no.~3, 351--399.

\bibitem{Zwi93}
B. Zwiebach.
\textit{Closed string field theory: quantum action and the Batalin--Vilkovisky master equation}.
Nuclear Physics B \textbf{390} (1993), no.~1, 33--152.

\bibitem{Deligne71}
P. Deligne.
\textit{Th\'eorie de Hodge: II}.
Inst. Hautes \'Etudes Sci. Publ. Math. No. 40 (1971), 5--57.

\bibitem{KashiwaraSchapira}
M. Kashiwara and P. Schapira.
\textit{Sheaves on Manifolds}.
Grundlehren der mathematischen Wissenschaften, vol. 292, Springer-Verlag, Berlin, 1990.

\bibitem{BlumPlauschin}
R. Blumenhagen and E. Plauschinn.
\textit{Introduction to Conformal Field Theory: With Applications to String Theory}.
Lecture Notes in Physics, vol. 779, Springer-Verlag, Berlin, 2009.

\bibitem{PP05}
A. Polishchuk and L. Positselski.
\textit{Quadratic Algebras}.
University Lecture Series \textbf{37}, American Mathematical Society, Providence, RI, 2005.

\bibitem{DI97}
J. Ding and K. Iohara.
\textit{Generalization of Drinfeld quantum affine algebras}.
Letters in Mathematical Physics \textbf{41} (1997), no.~2, 181--193.

\bibitem{BK86}
A. A. Belavin and V. G. Knizhnik.
\textit{Algebraic geometry and the geometry of quantum strings}.
Physics Letters B \textbf{168} (1986), no.~3, 201--206.

\bibitem{Bor92}
R. E. Borcherds.
\textit{Monstrous moonshine and monstrous Lie superalgebras}.
Inventiones Mathematicae \textbf{109} (1992), no.~1, 405--444.

\bibitem{Seg04}
G. Segal.
\textit{The definition of conformal field theory}.
In: \textit{Topology, Geometry and Quantum Field Theory}, London Math. Soc. Lecture Note Ser. \textbf{308}, Cambridge Univ. Press, 2004, pp.~421--577.

\bibitem{BV85}
D. Barbasch and D. Vogan.
\textit{Unipotent representations of complex semisimple groups}.
Annals of Mathematics \textbf{121} (1985), no.~1, 41--110.

\bibitem{Ger64}
M. Gerstenhaber.
\textit{On the deformation of rings and algebras}.
Annals of Mathematics \textbf{79} (1964), no.~1, 59--103.

\bibitem{Fuks86}
D. B. Fuks.
\textit{Cohomology of Infinite-Dimensional Lie Algebras}.
Consultants Bureau, New York, 1986.

\bibitem{Har77}
R. Hartshorne.
\textit{Algebraic Geometry}.
Graduate Texts in Mathematics \textbf{52}, Springer-Verlag, New York, 1977.

\bibitem{KO68}
N. M. Katz and T. Oda.
\textit{On the differentiation of De Rham cohomology classes with respect to parameters}.
J. Math. Kyoto Univ. \textbf{8} (1968), 199--213.

\bibitem{Ara74}
S. Ju. Arakelov.
\textit{An intersection theory for divisors on an arithmetic surface}.
Izv. Akad. Nauk SSSR Ser. Mat. \textbf{38} (1974), 1179--1192.

\bibitem{Fal84}
G. Faltings.
\textit{Calculus on arithmetic surfaces}.
Annals of Mathematics \textbf{119} (1984), no.~2, 387--424.

% ================================================================
% SESSION 56: REFERENCES FROM references/ DIRECTORY
% ================================================================

\bibitem{Mok25}
Siao Chi Mok.
\textit{Logarithmic Fulton--MacPherson configuration spaces}.
arXiv:2503.17563, 2025.
Note: Constructs a logarithmic analogue $\mathrm{FM}_n(X|D)$ of the
Fulton--MacPherson compactification for non-proper varieties $X \setminus D$,
with degeneration formulas for families of FM spaces.

\bibitem{MNO96}
Alexander Molev, Maxim Nazarov, and Grigori Olshanskii.
\textit{Yangians and classical Lie algebras}.
Russian Math.\ Surveys \textbf{51}(2):205--282, 1996.
Note: Systematic development of Yangians $Y(\mathfrak{g})$ for classical
Lie algebras via the RTT presentation.  Generators, relations, PBW theorem,
and the relationship to universal enveloping algebras.

\bibitem{Latyntsev23}
Alexei Latyntsev.
\textit{Factorisation quantum groups}.
arXiv:2312.07274, 2023.
Note: Develops vertex and factorization algebra analogues of
quasitriangular bialgebras, controlled by spectral $R$-matrices.
Appendix A.1 recollects chiral Koszul duality.

\bibitem{KhanZeng25}
Ahsan Z. Khan and Keyou Zeng.
\textit{Poisson vertex algebras and three dimensional gauge theory}.
arXiv:2502.13227, 2025.
Note: Mixed holomorphic-topological gauge theory in 3d from
freely generated Poisson vertex algebras;
$\lambda$-bracket Jacobi identity $\Leftrightarrow$ gauge invariance.

\bibitem{Zeng23}
Keyou Zeng.
\textit{Twisted holography and celestial holography from boundary chiral algebra}.
arXiv:2302.06693, 2023.
Note: Kaluza--Klein reduction of 6d holomorphic theories; boundary chiral
algebras as Koszul duals of bulk local operators; $A_\infty$ structure on
tangential Cauchy--Riemann cohomology via homotopy transfer.

\bibitem{CDG20}
Kevin Costello, Tudor Dimofte, and Davide Gaiotto.
\textit{Boundary chiral algebras and holomorphic twists}.
arXiv:2005.00083, 2020.
Note: Holomorphic twist of 3d $\mathcal{N}=2$ gauge theories;
bulk operators form commutative chiral algebra with Poisson bracket;
boundary operators are modules, with $A_\infty$ structures.

\bibitem{GKW24}
Davide Gaiotto, Justin Kulp, and Jingxiang Wu.
\textit{Higher operations in perturbation theory}.
arXiv:2403.13049, 2024.
Note: Higher operations from Feynman diagrams in holomorphic-topological
theories; quadratic axioms as Wess--Zumino consistency;
higher-dimensional formality theorem.

\bibitem{DNP25}
Tudor Dimofte, Wenjun Niu, and Victor Py.
\textit{Line operators in 3d holomorphic QFT: meromorphic tensor categories
and dg-shifted Yangians}.
arXiv:2508.11749, 2025.
Note: Lines are modules for the $A_\infty$-algebra $\mathcal{A}^!$
Koszul dual to bulk local operators; the OPE of lines is controlled
by a ``dg-shifted Yangian.''

\bibitem{Tamarkin00}
Dmitry E. Tamarkin.
\textit{The deformation complex of a $d$-algebra is a $(d+1)$-algebra}.
arXiv:math/0010072, 2000.
Note: Purely algebraic proof that the deformation complex of an $E_d$-algebra
is an $E_{d+1}$-algebra; Koszul duality for the $E_d$ operad.

\bibitem{DTT09}
V. A. Dolgushev, D. E. Tamarkin, and B. L. Tsygan.
\textit{Formality theorems for Hochschild complexes and their applications}.
arXiv:0901.0069, 2009.
Note: Formality of Hochschild cochain and chain complexes; applications
to deformation quantization.

\bibitem{CattaneoFelder99}
Alberto S. Cattaneo and Giovanni Felder.
\textit{A path integral approach to the Kontsevich quantization formula}.
Communications in Mathematical Physics \textbf{212} (2000), 591--611.
arXiv:math/9902090.
Note: QFT interpretation of Kontsevich's deformation quantization
via topological sigma model on the disk; BV quantization yields
the star product.

\bibitem{KontsevichSoibelman}
Maxim Kontsevich and Yan Soibelman.
\textit{Deformation Theory.\ I}.
Unpublished manuscript.
Note: Foundational treatment of deformation theory via $L_\infty$-algebras,
Maurer--Cartan moduli, and formal moduli problems.

\bibitem{BarNatan95}
D. Bar-Natan.
\textit{On the Vassiliev knot invariants}.
Topology \textbf{34} (1995), no. 2, 423--472.

\bibitem{Mirzakhani}
M. Mirzakhani.
\textit{Weil--Petersson volumes and intersection numbers of moduli spaces of bordered Riemann surfaces}.
J. Amer. Math. Soc. \textbf{20} (2007), no.~1, 1--23.
Note: Recursive formula for Weil--Petersson volumes of moduli spaces;
implies $\mathrm{Vol}(\overline{\mathcal{M}}_g) \sim (2g)!$ growth.

\bibitem{Quillen85}
D. Quillen.
\textit{Determinants of Cauchy--Riemann operators over a Riemann surface}.
Funktsional. Anal. i Prilozhen. \textbf{19} (1985), no.~1, 37--41.

\bibitem{Witten91}
E. Witten.
\textit{On quantum gauge theories in two dimensions}.
Comm. Math. Phys. \textbf{141} (1991), no.~1, 153--209.
Note: Computes symplectic volumes of moduli spaces of flat connections
via two-dimensional gauge theory; Verlinde formula as large-level
asymptotics.

\bibitem{Andersen92}
H. H. Andersen.
\textit{Tensor products of quantized tilting modules}.
Communications in Mathematical Physics \textbf{149} (1992), no.~1, 149--159.

\bibitem{Soergel97}
W. Soergel.
\textit{Kazhdan--Lusztig polynomials and a combinatoric for tilting modules}.
Representation Theory \textbf{1} (1997), 83--114.

\bibitem{Beauville}
A. Beauville.
\textit{Conformal blocks, fusion rules and the Verlinde formula}.
In: \textit{Proceedings of the Hirzebruch 65 Conference on Algebraic Geometry},
Israel Math. Conf. Proc. \textbf{9}, Bar-Ilan Univ., 1996, pp.~75--96.
arXiv:alg-geom/9405001.

\bibitem{NakanishiTsuchiya}
T. Nakanishi and A. Tsuchiya.
\textit{Level-rank duality of WZW models in conformal field theory}.
Communications in Mathematical Physics \textbf{144} (1992), no.~2, 351--372.

\bibitem{GKO85}
P. Goddard, A. Kent, and D. Olive.
\textit{Virasoro algebras and coset space models}.
Physics Letters B \textbf{152} (1985), no.~1--2, 88--92.

\bibitem{Kostant78}
B. Kostant.
\textit{On Whittaker vectors and representation theory}.
Inventiones Mathematicae \textbf{48} (1978), no.~2, 101--184.

\bibitem{HarishChandra}
Harish-Chandra.
\textit{On some applications of the universal enveloping algebra of a semisimple Lie algebra}.
Transactions of the American Mathematical Society \textbf{70} (1951), 28--96.

% === Programme I: Langlands ===

\bibitem{FT06}
Edward Frenkel and Constantin Teleman.
\textit{Self-extensions of Verma modules and differential forms on opers}.
Compositio Mathematica \textbf{142} (2006), no.~2, 477--500.
Note: Proves $\mathrm{Ext}^n(V_{\mathrm{crit}}, V_{\mathrm{crit}}) = \Omega^n(\mathrm{Op}_{\check{\mathfrak{g}}})$ for all~$n$,
identifying the full derived endomorphism algebra of the vacuum module at critical level
with differential forms on the space of opers.

\bibitem{Raskin24}
Sam Raskin.
\textit{Homological methods in semi-infinite contexts}.
arXiv:2002.01395v3, 2024.
Note: Establishes the \emph{Fundamental Local Equivalence} (FLE):
$\mathrm{KL}_{\mathrm{crit}}(\widehat{\mathfrak{g}}) \simeq \mathrm{Whit}(\widehat{\mathfrak{g}}^\vee)$,
an equivalence between the Kazhdan--Lusztig category at critical level and the
Whittaker category for the Langlands dual.  This is a cornerstone of the
geometric Langlands programme.

\bibitem{GLC24}
Dennis Gaitsgory et~al.
\textit{Proof of the geometric Langlands conjecture}.
arXiv:2405.03599, 2024.
Note: Completes the proof of the geometric Langlands conjecture for $\mathrm{GL}_n$
and more generally for reductive groups, building on the FLE of Raskin
and the categorical framework of Arinkin--Gaitsgory.

% === Programme II: N-complexes ===

\bibitem{Kapranov96}
Mikhail Kapranov.
\textit{On the $q$-analog of homological algebra}.
arXiv:q-alg/9611005, 1996.
Note: Introduces $N$-complexes ($d^N = 0$) as the natural algebraic
structure at roots of unity; homological algebra with $d^2 \neq 0$
but $d^N = 0$ for a prescribed~$N$.

\bibitem{KQ20}
Mikhail Khovanov and You Qi.
\textit{A faithful braid group action on the stable category of tricomplexes}.
SIGMA \textbf{16} (2020), 019.
arXiv:1911.04515.
Note: Develops the homotopy theory of $N$-complexes ($d^N = 0$);
constructs faithful braid group actions on derived categories of
$N$-complexes, connecting to quantum group categorification at
roots of unity.

% === Programme III: Monoidal coalgebras ===

\bibitem{BLPW16} T. Braden, A. Licata, N. Proudfoot and B. Webster, \emph{Quantizations of conical symplectic resolutions II: category $\mathcal{O}$ and symplectic duality}, Ast\'erisque \textbf{384} (2016), 75--179.

\bibitem{BL24}
Jack Booth and Andrey Lazarev.
\textit{Monoidal model structures on comodule categories}.
arXiv:2406.04684, 2024.
Note: Constructs monoidal model structures on categories of
coalgebras/comodules, providing the homotopical infrastructure
for bar-cobar adjunctions to be Quillen equivalences in the
monoidal setting.

% === Programme IV: E_n higher-dimensional ===

\bibitem{Idrissi22}
Najib Idrissi.
\textit{Real Homotopy of Configuration Spaces: Peccot Lecture,
Coll\`ege de France, March and May 2020}.
Lecture Notes in Mathematics, vol.~2303, Springer, 2022.
Note: Proves that Poincar\'e duality on configuration spaces
ensures $d^2 = 0$ for the propagator-based differential on
$\mathsf{E}_n$-algebras; provides explicit real models for
configuration space cohomology.

\bibitem{Knudsen18}
Ben Knudsen.
\textit{Higher enveloping algebras}.
Geometry \& Topology \textbf{22} (2018), no.~7, 4013--4066.
arXiv:1605.01391.
Note: Defines higher enveloping algebras $U_n(\mathfrak{g})$ as
the $\mathsf{E}_n$-analogue of the universal enveloping algebra;
proves PBW and Koszul duality theorems in the $\mathsf{E}_n$ setting.

% === Programme V: Vassiliev / Chern-Simons ===

\bibitem{CattaneoMnev10}
Alberto~S. Cattaneo and Pavel Mnev.
\textit{Remarks on Chern--Simons invariants}.
Communications in Mathematical Physics \textbf{293} (2010), no.~3, 803--836.
arXiv:0811.2045.
Note: Establishes that the bridge between holomorphic and topological
Chern--Simons gauge equivalence classes provides the correct comparison;
configuration space integrals on $3$-manifolds produce Vassiliev invariants
via BV integration.

% === Programme VI-a: BRST chain map and BV quantization ===

\bibitem{SiLi12}
Si Li.
\textit{Feynman graph integrals and almost modular forms}.
Communications in Mathematical Physics \textbf{314} (2012), no.~2, 299--341.
arXiv:1112.4015.
Note: Establishes BV quantization for holomorphic Chern--Simons theory
via Feynman graph integrals; produces the $A_\infty$ structure on the
deformation complex from perturbative path integrals.  The genus-$1$
contributions are expressed as almost-modular forms.

\bibitem{Arkhipov97}
Sergey Arkhipov.
\textit{Semi-infinite cohomology of associative algebras and bar-duality}.
International Mathematics Research Notices \textbf{1997}, no.~17, 833--863.
Note: Constructs a genus-$0$ chain map between the semi-infinite (BRST)
complex and the bar complex for associative algebras; the chiral
generalization of this map is one component of the bar-BRST identification.

\bibitem{Markl23}
Martin Markl.
\textit{On the origin of higher braces and higher-order derivations}.
Journal of Physics A \textbf{56} (2023), 205202.
arXiv:2211.08197.
Note: Clarifies the relationship between BV brackets and bar complex
structures via higher-order derivations; relevant to the chain-level
identification of BV data with bar complex operations.

\bibitem{SiLi16}
Si Li.
\textit{Vertex algebras and quantum master equation}.
Journal of Differential Geometry \textbf{123} (2023), no.~1, 127--179.
arXiv:1612.01292.
Note: Establishes BV quantization for holomorphic field theories on
Riemann surfaces, producing $A_\infty$ structures from Feynman
diagram integrals; the chain-level output is the curved bar complex.

% === Programme VII: NC Hodge / weight filtrations ===

\bibitem{BG24}
Christin Bibby and Nir Gadish.
\textit{Combinatorics of orbit configuration spaces}.
S\'eminaire Lotharingien de Combinatoire \textbf{91B} (2024), Art.~23.
Note: Proves that the weight filtration on configuration spaces
$C_n(X)$ for complements of hyperplane arrangements is explicitly
computable via combinatorial data (intersection lattice); provides
effective algorithms for computing the mixed Hodge structure.

% === Programme VIII-a: W-algebra orbits ===

\bibitem{AvE23}
Tomoyuki Arakawa and Jethro van Ekeren.
\textit{Rationality and fusion rules for $\mathcal{W}$-algebras with
hook-type nilpotent}.
arXiv:2306.08112, 2023.
Note: Proves rationality and explicit Koszul duality for hook-type
W-algebras $\mathcal{W}^k(\mathfrak{sl}_N, f_{\mathrm{hook}})$;
extends the Feigin--Frenkel program beyond the principal nilpotent.

% BV85 already defined above (line ~1033); note on BV duality added there

\bibitem{Achar03}
Pramod~N. Achar.
\textit{An order-reversing duality map for conjugacy classes in
Lusztig's canonical quotient}.
Transformation Groups \textbf{8} (2003), no.~4, 303--323.
Note: Studies properties of BV duality including its relation
to Spaltenstein duality and the closure ordering on nilpotent orbits.

\bibitem{Knu83} F. F. Knudsen, \emph{The projectivity of the moduli space of stable curves, II: The stacks $M_{g,n}$}, Math. Scand. \textbf{52} (1983), 161--199.

\bibitem{FGZ86} I. B. Frenkel, H. Garland, and G. J. Zuckerman, \emph{Semi-infinite cohomology and string theory}, Proc. Natl. Acad. Sci. USA \textbf{83} (1986), no. 22, 8442--8446.

\bibitem{EM66} S. Eilenberg and J. C. Moore, \emph{Homology and fibrations. I. Coalgebras, cotensor product and its derived functors}, Comment. Math. Helv. \textbf{40} (1966), 199--236.

\bibitem{Wei94} C. A. Weibel, \emph{An Introduction to Homological Algebra}, Cambridge Studies in Advanced Mathematics, vol. 38, Cambridge University Press, Cambridge, 1994.

\bibitem{BoSch93} P. Bouwknegt and K. Schoutens, \emph{$\mathcal{W}$-symmetry in conformal field theory}, Phys. Rep. \textbf{223} (1993), no. 4, 183--276.

\bibitem{PTVV13} T. Pantev, B. To\"en, M. Vaqui\'e, and G. Vezzosi, \emph{Shifted symplectic structures}, Publ. Math. Inst. Hautes \'Etudes Sci. \textbf{117} (2013), 271--328.

\bibitem{CPTVV17} D. Calaque, T. Pantev, B. To\"en, M. Vaqui\'e, and G. Vezzosi, \emph{Shifted Poisson structures and deformation quantization}, J. Topol. \textbf{10} (2017), no. 2, 483--584.

\bibitem{Pridham17} J. P. Pridham, \emph{Shifted Poisson and symplectic structures on derived $N$-stacks}, J. Topol. \textbf{10} (2017), no. 1, 178--210.

\end{thebibliography}
